% LSPV-4 REPORT::WRAPINV1::2026-10-03T06:45:00-07:00
\documentclass[11pt]{article}
\usepackage[margin=.85in]{geometry}
\usepackage{amsmath,amssymb}
\begin{document}
\section*{LSPV-4 WRAPINV1: universal arbitrary-algorithm transfer requires wrapper invariance}
\textbf{STATUS: OPEN.} P versus NP remains open. This report proves an exact logical guard on any universal arbitrary-SAT-algorithm transfer. It proves no SAT lower bound, no SU/vMS asymptotic theorem, and no nonrelativizing or nonalgebrizing hinge.

\section*{CORPUS REVIEWED}
Before substantive work this run fetched and obeyed AGENTS.md, CODEX.md, and the repository pnp-research skill; surveyed root main.tex, current tasks TASK-001 through TASK-007, verified MAIN 0058, governance surfaces, and the active HEC/CRS-PPC/CCSPC/RCCN/Search-G branches. The immediate LSPV-3 predecessor BOTTLENECK1 was read back with its expected report header, substantive body, document terminator, and verification sentinel. BOTTLENECK1 remains exponent accounting only and does not block this run. The verified LSPV-4 lineage through STATICNP1 was reconciled, with TRACECLOSE1 retained as the exact negative baseline.

Repository state contains no verified BLACKBOXCUT1 report and no exact preserved failed BLACKBOXCUT1 scientific body. It was therefore not reconstructed or retried. No other exact preserved failed LSPV-4 body was found in GitHub requiring recovery.

Primary literature was refreshed around relativization, algebrization, Natural Proofs, black-box lower bounds, and meta-complexity. Baker--Gill--Solovay remains the exact relativization barrier. Razborov--Rudich remains the Natural Proofs guard for broad constructive circuit properties under its hypotheses. Aaronson--Wigderson remains the algebrization guard. Recent meta-mathematics work of Garlik et al.\ studies proof-theoretic strength of algebraic lower bounds, but does not supply an arbitrary-SAT-algorithm acquisition theorem. Black-box query lower bounds such as Juma--Kabanets--Rackoff--Shpilka are exact only in their oracle model and therefore serve as model-boundary controls, not unrestricted SAT transfers.

\section*{Target and quantifiers}
The current load-bearing target has the universal form
\[
\forall A\;\Bigl[
A\text{ is a deterministic polynomial-time decider for SAT}
\Longrightarrow
\exists F\;{\cal R}(A,F)
\Bigr],
\]
where \({\cal R}(A,F)\) is intended to force SAME\_REALIZATION acquisition, provenance, reconstruction, or a no-sharing resource.

Let \({\cal W}\) be any family of semantics-preserving polynomial-overhead wrappers. For \(W\in{\cal W}\), \(W(A)\) decides exactly the same SAT language as \(A\) and remains polynomial time. DEAD, COPY, DELAY by polynomial factors, INTERP, and polynomial RECOMP transformations are admissible examples when their declared implementation preserves correctness and polynomial time.

\section*{WRAPINV1 theorem}
\textbf{Theorem.} Suppose
\[
\forall A\in{\cal A}_{SAT}^{P}\;\exists F\;{\cal R}(A,F)
\tag{1}
\]
is claimed, where \({\cal A}_{SAT}^{P}\) is the class of all deterministic polynomial-time SAT deciders. Let \(W\) be any semantics-preserving polynomial-overhead transformation satisfying
\[
A\in{\cal A}_{SAT}^{P}\Longrightarrow W(A)\in{\cal A}_{SAT}^{P}.
\tag{2}
\]
Then the universal claim necessarily implies
\[
\forall A\in{\cal A}_{SAT}^{P}\;\exists F'\;{\cal R}(W(A),F').
\tag{3}
\]
Consequently, if there exists an \(A\in{\cal A}_{SAT}^{P}\) and an admissible wrapper \(W\) for which
\[
\forall F'\;\neg{\cal R}(W(A),F'),
\tag{4}
\]
then the universal transfer (1) is false.

\paragraph{Proof.}
By (2), \(W(A)\) is itself an element of \({\cal A}_{SAT}^{P}\). Instantiating the universal quantifier in (1) with \(W(A)\) gives (3). Therefore (4) contradicts (1). \(\square\)

This is elementary, but it closes an important quantifier loophole: a resource attached to one implementation of a hypothetical solver is not universally forced unless the transfer survives every admissible polynomial semantics-preserving implementation change, or the theorem supplies a representation-invariant replacement resource for the wrapped implementation.

\section*{Corollary: code-identity and exact-history resources}
Let \(P(A)\) be an implementation property used as a necessary premise for the claimed lower-bound resource. If there exists an admissible wrapper \(W\) and a polynomial SAT decider \(A\) such that
\[
P(A)\quad\text{but}\quad \neg P(W(A)),
\]
then \(P\) cannot be asserted as a necessary property of every polynomial SAT decider merely from SAT correctness and polynomial runtime.

In particular, exact instruction sequence, exact state naming, exact trace length, exact query chronology when chronology is wrapper-mutable, exact memory address identity, exact provenance labels, and exact designated intermediate objects are not universal semantic necessities without an additional invariant theorem.

This does not say that all implementation-sensitive information is useless. It says that a successful non-black-box route must identify a property that is forced for every polynomial implementation after the quantifier is instantiated, or prove a canonical representation theorem whose acquisition and conversion costs are fully charged and themselves wrapper-safe.

\section*{SAME\_REALIZATION consequence}
SAME\_REALIZATION remains necessary when several charged resources are combined, but SAME\_REALIZATION does not itself provide wrapper invariance. The valid quantifier shape is
\[
\forall A\;\exists F_A\;\exists Q_A\;
\Bigl[
\operatorname{SameRealization}(A,F_A,Q_A)
\wedge {\cal C}(A,F_A,Q_A)
\Bigr].
\]
One may not prove \({\cal C}\) for a convenient implementation \(A_0\), alter \(A_0\) with DEAD/COPY/DELAY/INTERP/RECOMP, and continue charging the original \(F_{A_0},Q_{A_0}\) unless an explicit transport theorem preserves the resource and charges the transport. Nor may one invert the quantifiers to a single fixed \(F\) or \(Q\) shared by all algorithms.

\section*{Hostile audit of active branches}
\paragraph{HEC/CRS-PPC/CCSPC.}
Any static-to-causal compiler must survive an implementation that recomputes a valid current response instead of retaining the reference transcript, interprets another machine, inserts dead computation, duplicates internal work, or polynomially delays epochs. If the claimed hard causal cut depends on the reference implementation's epoch boundary rather than a wrapper-invariant input property, WRAPINV1 blocks universal promotion.

\paragraph{RCCN.}
A routed-store lower bound can remain exact for its declared network. But if an equivalent polynomial implementation flattens the routing, changes storage locations, or recomputes messages, the original routed decomposition is not a universal resource. The bridge must quantify over the wrapped implementation and prove a representation-invariant inequality.

\paragraph{Search-G and PRIME-LQP.}
Precommitment, prime tags, hashes, and provenance histories are not universally forced merely because one implementation emits them. RECOMP can choose a different valid witness when allowed; INTERP changes realization; DEAD/COPY/DELAY alter trace provenance. A surviving Search-G theorem must define the audit resource extensionally enough to survive these transformations while remaining SAME\_REALIZATION and fully charged.

\paragraph{SU/vMS/Hadamard.}
WRAPINV1 does not attack the fixed algebraic identities themselves. It attacks the transfer from those identities to arbitrary computation. Even if the exact source-compatible characteristic-two rank grows superlinearly, a separate theorem must show that every polynomial SAT decider, including wrapped extensionally equivalent deciders, acquires a representation-invariant resource tied to that rank.

\section*{Easy-family controls}
Horn, dual-Horn, 2-SAT, affine/Tseitin, bounded-treewidth, padding, renaming, auxiliary-variable, and independent-product controls remain mandatory. They sharpen WRAPINV1: if a proposed implementation property can be manufactured by a wrapper around an easy-family solver, then the property is not by itself evidence of semantic hardness. Conversely, if a proposed invariant falsely charges an easy family after harmless wrapping, the invariant or transfer is mis-specified.

\section*{Barrier audit}
WRAPINV1 is a logical closure guard and does not cross relativization. The same quantifier argument works for oracle machines whenever the wrapper and algorithm receive the same oracle. Baker--Gill--Solovay therefore remains fully applicable to an otherwise relativizing proof spine.

WRAPINV1 also supplies no nonalgebrizing mechanism. Aaronson--Wigderson's algebrization barrier remains a separate audit. A future implementation-sensitive cut that is genuinely nonrelativizing must still be tested after low-degree-extension oracle access.

Natural Proofs is not directly triggered by WRAPINV1 because no large constructive property of Boolean functions is asserted. If the proposed wrapper-invariant replacement is promoted to an efficiently recognizable large property useful against general circuits, the Razborov--Rudich audit becomes mandatory.

\section*{Cross-domain and PQ classification}
Black-box lower bounds from optimization, communication, cryptography, or polynomial-oracle query models are \textbf{FAILED TRANSFER} when imported directly as unrestricted SAT runtime lower bounds; their restricted interface is load-bearing. LWE-style worst-case/average-case black-box extraction remains an \textbf{ANALOGY/CANDIDATE}: an external extractor could be valuable only after a theorem shows that every SAT solver's black-box behavior yields the source-matched object with fully charged polynomial overhead. Schur-square and quadratic-relation tools remain \textbf{EXACT OBJECT/THEOREM TRANSFER} only for the matching code-algebra object, not for arbitrary-algorithm acquisition.

\section*{DUPLICATION / CONTRADICTION AUDIT}
The corpus already mandates DEAD/COPY/DELAY/INTERP/RECOMP and warns that representation size is not runtime. WRAPINV1 does not merely repeat that heuristic. It formalizes the exact universal-quantifier reason these wrappers are mandatory: closure of the quantified algorithm class under an admissible wrapper forces the transfer theorem to hold again for the wrapped algorithm.

TRACECLOSE1 remains distinct: it proves polynomial reacquisition of deterministic polynomial-trace functionals. STATICNP1 remains distinct: it closes NP-level continuation search under polynomial SAT staticization and \(P=NP\). BOTTLENECK1 remains accounting. No contradiction with the structural PAIRDEG1/FULLPAIR6 line was found.

\section*{WHAT CHANGED RELATIVE TO ROOT MAIN}
Root main.tex lists wrappers as hostile controls and leaves representation-invariant no-sharing OPEN. WRAPINV1 upgrades the wrapper requirement from methodological hygiene to an exact necessary condition on any universal arbitrary-algorithm transfer. It does not close representation-invariant no-sharing; it specifies a falsification criterion for candidate transfers.

\section*{ADOPTED / REJECTED / RETIRED PRIOR COMPONENTS}
\textbf{Adopted:} TRACECLOSE1 as the negative baseline; LEGALFLAT1, RESPSEL1, GLOBALGAME1, STATICNP1; fixed-realization SU/vMS algebra within its verified scope; the restricted HEC/CRS-PPC/CCSPC/RCCN/Search-G theorems within their declared models; and DEAD/COPY/DELAY/INTERP/RECOMP as admissible hostile controls when polynomial correctness is preserved.

\textbf{Rejected:} treating source dependence, exact code identity, exact transcript identity, exact memory addresses, designated provenance tags, or a chosen modular decomposition as universally necessary without a wrapper-invariant transfer theorem.

\textbf{FALSIFIED/RETIRED as generic hinges:} any arbitrary-algorithm transfer whose required property is destroyed by an admissible polynomial semantics-preserving wrapper and has no replacement resource satisfying the theorem for the wrapped algorithm.

\section*{Components consistent with \(P=NP\)}
TRACECLOSE1, LEGALFLAT1, RESPSEL1, GLOBALGAME1, STATICNP1, WRAPINV1, fixed-realization SU/vMS algebra, wrapper-invariant quantifier hygiene, and current reconstruction accounting are all consistent with \(P=NP\). None supplies the final contradiction.

\section*{One next falsification target}
Take the strongest current CRS-E/Search-G or SU/vMS acquisition statement and write its resource predicate explicitly as
\[
{\cal R}(A,F,Q).
\]
For each wrapper DEAD, COPY, DELAY, INTERP, and RECOMP, prove
\[
\forall A\in{\cal A}_{SAT}^{P}\;\exists F',Q'\;
{\cal R}(W(A),F',Q')
\]
with all transport/acquisition costs charged, or exhibit a wrapper that destroys the predicate. The first destroyed predicate is retired as a universal hinge. Only a predicate surviving the full wrapper battery should proceed to maximal semantic-oracle and algebraic-oracle tests.

\section*{Strict status}
WRAPINV1 is \textbf{PROVED/EXACT} as the stated universal-quantifier closure theorem. Application to a concrete acquisition/no-sharing predicate is \textbf{CONDITIONAL} until that predicate and wrapper maps are source-transcribed. Exact SU/vMS asymptotic growth, arbitrary-algorithm acquisition, representation-invariant no-sharing, fully charged reconstruction, and a genuinely nonrelativizing and nonalgebrizing contradiction remain \textbf{OPEN}. P versus NP remains open.

\section*{References}
T. Baker, J. Gill, and R. Solovay, Relativizations of the P=?NP Question, SIAM Journal on Computing 4(4):431--442, 1975.

A. A. Razborov and S. Rudich, Natural Proofs, Journal of Computer and System Sciences 55(1):24--35, 1997.

S. Aaronson and A. Wigderson, Algebrization: A New Barrier in Complexity Theory, ACM Transactions on Computation Theory 1(1), 2009.

A. Juma, V. Kabanets, C. Rackoff, and A. Shpilka, The Black-Box Query Complexity of Polynomial Summation, Computational Complexity 18(1):59--79, 2009.

M. Garlik, S. Gryaznov, J. Lu, R. Santhanam, and I. Tzameret, Meta-Mathematics of Algebraic Complexity, LICS 2026, LIPIcs 380:49:1--49:25, 2026.
\end{document}
% END_OF_VERIFIED_REPORT::LSPV-4::WRAPINV1
